Le système d'acquisition de l'ITIPhone intègre quatre capteurs : un module GNSS
(DFRobot TEL0157), une centrale inertielle (MiniIMU-9 v6), un LiDAR 2D
(RPLIDAR A1M8) et une caméra (Raspberry Pi Camera Module 3). Cette section
justifie le choix des capteurs retenus pour la fusion de données.

\subsubsection{Module GNSS}

Le module GNSS fournit une position absolue en latitude et longitude à une
fréquence comprise entre 1 et 10~Hz. Son erreur est de nature
\textbf{stochastique et bornée} : les mesures oscillent aléatoirement autour
de la vraie position, sans dérive temporelle. Ce comportement est modélisé
par un bruit blanc gaussien de variance $\sigma^2_{\text{GPS}}$, typiquement de
l'ordre de quelques mètres carrés en milieu urbain.

Cependant, le GPS seul présente deux limitations majeures pour une estimation
de trajectoire robuste :

\begin{itemize}
    \item \textbf{Fréquence d'acquisition faible} : une mesure toutes les
    100~ms à 1~s est insuffisante pour restituer fidèlement une trajectoire
    lors d'un déplacement dynamique. Entre deux mesures, la position est
    inconnue.
    \item \textbf{Disponibilité conditionnelle} : en environnement urbain
    dense ou sous couvert végétal, le signal peut être temporairement
    indisponible, rendant l'estimation impossible sans capteur de substitution.
\end{itemize}

\subsubsection{Centrale inertielle IMU}

L'IMU MiniIMU-9 v6 intègre un accéléromètre, un gyroscope et un magnétomètre,
opérant à une fréquence typique de 100~Hz. La position est estimée par double
intégration numérique de l'accélération mesurée dans le repère monde :

\begin{align}
    \mathbf{v}(t) &= \mathbf{v}(t-1) + \mathbf{a}_{\text{monde}}(t) \cdot \Delta t \\
    \mathbf{x}(t) &= \mathbf{x}(t-1) + \mathbf{v}(t) \cdot \Delta t
\end{align}

L'IMU présente une limitation fondamentale : son erreur est \textbf{non
bornée}. Le biais $b$ de l'accéléromètre, après double intégration, produit
une erreur de position croissant en $t^2$ :

\begin{equation}
    \epsilon_{\text{position}}(t) \approx \frac{1}{2} \cdot b \cdot t^2
\end{equation}

Pour un biais typique de $b = 0{,}01$~m/s², l'erreur atteint 18~mètres après
60 secondes. \textbf{L'IMU seul est donc inutilisable pour une navigation de
longue durée.}

\subsubsection{LiDAR 2D et caméra}

Ces deux capteurs peuvent théoriquement fournir une estimation de déplacement
relatif (via l'algorithme ICP pour le LiDAR, ou l'odométrie visuelle pour la
caméra). Ils ont été écartés pour la fusion pour trois raisons :

\begin{itemize}
    \item Ils ne fournissent qu'un \textbf{déplacement relatif}, sans
    référence de position absolue — fonctionnellement équivalent à l'IMU,
    mais bien plus complexe à implémenter.
    \item Leurs performances en extérieur sont dégradées : l'ICP diverge dans
    les environnements peu structurés, et l'odométrie visuelle est sensible
     aux conditions lumineuses.
    \item Ces algorithmes sont coûteux en calcul pour une Raspberry Pi 4
    embarquée.
\end{itemize}

Ces capteurs seront valorisés dans le cadre de la Partie~2.

\subsection{Complémentarité GPS + IMU}

Le critère fondamental pour sélectionner les capteurs d'une fusion est la
\textbf{complémentarité de leurs caractéristiques d'erreur}. Le tableau
suivant résume les propriétés de chaque capteur :

\begin{table}[htbp]
    \centering
    \caption{Comparaison des caractéristiques d'erreur du GPS et de l'IMU}
    \begin{tabular}{lcc}
        \hline
        \textbf{Caractéristique} & \textbf{GPS} & \textbf{IMU} \\
        \hline
        Type d'erreur        & Bruit gaussien borné & Dérive en $t^2$ \\
        Fréquence            & 1--10~Hz             & $\sim$100~Hz    \\
        Fiabilité long terme & Bonne                & Mauvaise        \\
        Fiabilité court terme & Mauvaise (bruit)    & Bonne           \\
        Disponibilité        & Conditionnelle       & Permanente      \\
        \hline
    \end{tabular}
    \label{tab:capteurs}
\end{table}

La combinaison GPS + IMU est complémentaire sur chacun des problèmes identifiés :

\begin{enumerate}
    \item \textbf{Fréquence basse du GPS} : l'IMU à 100~Hz interpole la
    trajectoire entre deux mesures GPS, produisant un tracé continu et
    physiquement cohérent.
    \item \textbf{Perte de signal GPS} : l'IMU prend le relais en navigation
    autonome. Le filtre de Kalman modélise la dégradation progressive de
    la confiance via la matrice de covariance $\mathbf{P}$.
    \item \textbf{Bruit du GPS} : le filtre de Kalman combine les deux
    sources de manière optimale. On peut montrer que la variance de
    l'estimation fusionnée vérifie : la fusion de deux capteurs bruités produit
    donc \textbf{toujours une estimation plus précise} que le meilleur capteur
    pris seul, à condition que leurs bruits soient décorrélés — ce qui est
    vérifié ici, les deux capteurs reposant sur des principes physiques distincts.
\end{enumerate}
